%% cover_letter.tex — standalone compile for MSSP submission
\documentclass[11pt]{article}
\usepackage[a4paper,margin=2.5cm]{geometry}
\usepackage{times}
\usepackage{graphicx}
\usepackage[hidelinks]{hyperref}
\pagestyle{empty}
\setlength{\parskip}{0.6em}

\begin{document}

\begin{flushright}
\today
\end{flushright}

\bigskip
Dear Editors-in-Chief of \textit{Mechanical Systems and Signal Processing},

\medskip
Please consider the enclosed manuscript, ``Constant predictions in
multi-condition aero-engine RUL models: causes, early detection, and
condition-adaptive network design'', for publication as a research paper.

Hybrid networks that combine convolutional, recurrent and attention layers
are common in work on aero-engine remaining useful life (RUL), and they are
usually assessed with one training run per configuration on NASA C-MAPSS.
We audited 161 runs of one such model and found that 37\% of FD002 runs and
24\% of FD004 runs end by printing one fixed number for every engine, the
mean of the training labels. None of the 81 single-condition runs did. A
formula predicts the test RMSE of these failed runs to within 0.01 cycles,
and the PHM08 score used in service is 15--250$\times$ worse than for a
working model. Because the training loss falls normally, a single RMSE near
43.8 looks like mediocre tuning; this is easy to miss in a one-number report.

Probes of saved checkpoints and six instrumented re-training runs explain
why. Every multi-condition run first spends 16--34 epochs outputting that
average, with its internal features recording the operating condition rather
than wear and its attention weighting every time step equally. Training
starts to improve again only when gradients in the convolutional branch
suddenly grow about a hundredfold, before early stopping ends the run. From
this we set window length and depth by the number of operating conditions,
keep several branches able to carry the degradation signal, and add three
cheap checks after training (prediction spread, attention entropy and a long
flat stretch in validation RMSE). Over five seeds the rules trained
reliably on FD001/FD002/FD003; on FD004 the seed still decides, with six of
eight runs failing and the two that worked agreeing at 19.53$\pm$0.07 RMSE.
The checks catch those failures early enough to start a fresh run.

We think the manuscript will interest MSSP readers because it concerns
whether learned predictors used in machinery health monitoring can be
trusted to train again, and because the per-epoch gradient, attention and
feature traces released with the paper give a concrete basis for studying
the optimisation. The same checks apply to other sensor-regression tasks
with several operating conditions. A Word file of supplementary material
accompanies the submission.

The manuscript is original, has not been published, and is not under
consideration elsewhere. All authors have approved it and declare no
competing interests. This work was supported by the Jiangxi Provincial
Early-Career Talent Fund (grant no. 20252BEJ730138) and the Jiangxi Provincial
Talent Introduction and Cultivation Fund (grant no. 117C00230023). The
processed data, experiment scripts, training logs and analysis code are
available in a tagged release of our repository at
\url{https://github.com/HugoZha/RUL-CAF/releases}. The C-MAPSS data are a
U.S.\ Government work in the public domain distributed by the NASA
Prognostics Data Repository; we acknowledge the repository and the data
donors and cite Saxena and Goebel (2008) as the source dataset.

Following the journal's practice, we would like to suggest the following
colleagues as potential reviewers. Each is active in an area directly
relevant to this work, as noted below. We have no recent collaboration or
other conflict of interest with any of them, and the suggestion is offered
only for the editors' consideration.

\begin{enumerate}
  \item Prof. Zhigang Tian, Department of Mechanical Engineering, University
  of Alberta, Edmonton, Canada. E-mail: ztian@ualberta.ca.
  \emph{Expertise in machine prognostics and remaining-useful-life prediction
  under varying operating conditions, including recent work in this journal.}

  \item Dr. Meng Ma, School of Mechanical Engineering, Xi'an Jiaotong
  University, Xi'an, China. E-mail: meng\_ma@xjtu.edu.cn.
  \emph{Expertise in data-driven diagnostics and health management of
  propulsion systems and in deep learning for condition monitoring.}

  \item Dr. Zhenghua Chen, James Watt School of Engineering, University of
  Glasgow, Glasgow, UK. E-mail: Zhenghua.Chen@glasgow.ac.uk.
  \emph{Expertise in turbofan-engine RUL prediction and time-series
  representation learning, including a recent review of the field.}

  \item Prof. Michael Pecht, Department of Mechanical Engineering, University
  of Maryland, College Park, USA. E-mail: pecht@umd.edu.
  \emph{Broad expertise in prognostics and health management and in
  data-driven and fusion-based approaches to remaining useful life.}

  \item Dr. Spyridon Plakias, Department of Electrical and Computer
  Engineering, Democritus University of Thrace, Xanthi, Greece.
  E-mail: splakias@ee.duth.gr.
  \emph{Expertise in deep learning for aircraft-engine RUL estimation on
  the NASA C-MAPSS data and in attention-based models.}
\end{enumerate}

Thank you for your consideration.

\medskip
Sincerely,\\[0.4em]

\noindent\includegraphics[width=4.2cm]{yuguo_zha_signature.png}\\[0.2em]
Yuguo Zha\\
Corresponding author\\
E-mail: hugozha.nchu@gmail.com

\end{document}
